% FAC-EN-U0045
% French source: fac_body.tex lines 1604--1633.
% Authority: PDF physical pp. 44--45; printed pp. 239--240 (no. 47 complete).
% Semantic entities: FAC-SEM-NO-II-3-47; FAC-SEM-PROP-II-3-8;
% FAC-SEM-THM-II-3-4; FAC-SEM-THM-II-3-5;
% FAC-SEM-COR-II-3-5-1; FAC-SEM-COR-II-3-5-2.

\medskip
\phantomsection\label{fac:II:3:no:47}
\textbf{47. Coverings of algebraic varieties by affine open subsets.}

\phantomsection\label{fac:II:3:prop:8}
\textsc{Proposition 8.} \textit{Let $X$ be an affine variety, and let
$\mathfrak{U}=\{U_i\}_{i\in I}$ be a finite covering of $X$ by affine open
subsets. If $\mathcal{F}$ is a coherent algebraic sheaf on $X$, then
$H^q(\mathfrak{U},\mathcal{F})=0$ for every $q>0$.}

By \hyperref[fac:II:3:no:42:prop:3]{Proposition~3}, there are regular functions $P_j$ on $X$ such that the
covering $\mathfrak{B}=\{X_{P_j}\}$ is finer than $\mathfrak{U}$. For every
$(i_0,\ldots,i_p)$, the covering $\mathfrak{B}_{i_0\ldots i_p}$ induced by
$\mathfrak{B}$ on $U_{i_0\ldots i_p}$ is defined by the restrictions of the
$P_j$ to $U_{i_0\ldots i_p}$. Since $U_{i_0\ldots i_p}$ is an affine variety
by \hyperref[fac:II:3:no:42:prop:1]{Proposition~1}, \hyperref[fac:II:3:no:46:thm:3]{Theorem~3} applies to it, and therefore
$H^q(\mathfrak{B}_{i_0\ldots i_p},\mathcal{F})=0$ for every $q>0$. Applying
\hyperref[fac:I:4:no:29:prop:5]{Proposition~5 of no.~29}, we then obtain
\[
H^q(\mathfrak{U},\mathcal{F})=H^q(\mathfrak{B},\mathcal{F});
\]
and since $H^q(\mathfrak{B},\mathcal{F})=0$ for $q>0$ by \hyperref[fac:II:3:no:46:thm:3]{Theorem~3}, the
proposition is proved.

\phantomsection\label{fac:II:3:thm:4}
\textsc{Theorem 4.} \textit{Let $X$ be an algebraic variety, let
$\mathcal{F}$ be a coherent algebraic sheaf on $X$, and let
$\mathfrak{U}=\{U_i\}_{i\in I}$ be a finite covering of $X$ by affine open
subsets. The homomorphism
$\sigma(\mathfrak{U}):H^n(\mathfrak{U},\mathcal{F})\to H^n(X,\mathcal{F})$
is bijective for every $n\geq0$.}

Consider the family $\mathfrak{B}^{\alpha}$ of finite coverings of $X$ by
affine open subsets. By \hyperref[fac:II:3:no:42:thm:1:cor]{the corollary to Theorem~1}, these coverings are
arbitrarily fine. Moreover, for every system $(i_0,\ldots,i_p)$, the covering
$\mathfrak{B}^{\alpha}_{i_0\ldots i_p}$ induced by
$\mathfrak{B}^{\alpha}$ on $U_{i_0\ldots i_p}$ is a covering by affine open
subsets, by \hyperref[fac:II:3:no:42:prop:1]{Proposition~1}. Hence \hyperref[fac:II:3:prop:8]{Proposition~8} gives
$H^q(\mathfrak{B}^{\alpha}_{i_0\ldots i_p},\mathcal{F})=0$ for $q>0$.
\hyperref[fac:I:4:no:29:thm:1:item:a]{Conditions (a)} and \hyperref[fac:I:4:no:29:thm:1:item:b]{(b)} of \hyperref[fac:I:4:no:29:thm:1]{Theorem~1 of no.~29} are therefore satisfied, and
the theorem follows.

\phantomsection\label{fac:II:3:thm:5}
\textsc{Theorem 5.} \textit{Let $X$ be an algebraic variety, and let
$\mathfrak{U}=\{U_i\}_{i\in I}$ be a finite covering of $X$ by affine open
subsets. Let
$0\to\mathcal{A}\to\mathcal{B}\to\mathcal{C}\to0$ be an exact sequence of
sheaves on $X$, where $\mathcal{A}$ is coherent algebraic. The canonical
homomorphism
$H_0^q(\mathfrak{U},\mathcal{C})\to H^q(\mathfrak{U},\mathcal{C})$
(cf. \hyperref[fac:I:3:no:24]{no.~24}) is bijective for every $q\geq0$.}

It is clearly enough to show that
$C_0(\mathfrak{U},\mathcal{C})=C(\mathfrak{U},\mathcal{C})$; that is, every
section of $\mathcal{C}$ over $U_{i_0\ldots i_q}$ is the image of a section
of $\mathcal{B}$ over $U_{i_0\ldots i_q}$. This follows from \hyperref[fac:II:3:no:46:thm:3:cor:2]{Corollary~2 to
Theorem~3}.

\phantomsection\label{fac:II:3:thm:5:cor:1}
\textsc{Corollary 1.} \textit{Let $X$ be an algebraic variety, and let
$0\to\mathcal{A}\to\mathcal{B}\to\mathcal{C}\to0$ be an exact sequence of
sheaves on $X$, where $\mathcal{A}$ is coherent algebraic. The canonical
homomorphism $H_0^q(X,\mathcal{C})\to H^q(X,\mathcal{C})$ is bijective for
every $q\geq0$.}

This is an immediate consequence of \hyperref[fac:II:3:no:42:thm:1]{Theorems~1} \hyperref[fac:II:3:thm:5]{and~5}.

\phantomsection\label{fac:II:3:thm:5:cor:2}
\textsc{Corollary 2.} \textit{There is an exact sequence}
\[
\cdots\longrightarrow H^q(X,\mathcal{B})
\longrightarrow H^q(X,\mathcal{C})
\longrightarrow H^{q+1}(X,\mathcal{A})
\longrightarrow H^{q+1}(X,\mathcal{B})
\longrightarrow\cdots .
\]
